\begin{table*}[t]
\centering
\small
\resizebox{\textwidth}{!}{%
\begin{tabular}{llll}
\toprule
\textbf{Metric} & \textbf{Evaluation type} & \textbf{What it probes} & \textbf{Reported score} \\
\midrule
Validation MLM loss & In-domain held-out MLM & Fit to the held-out BabyLM-style validation split under the training objective & Mean cross-entropy; lower is better \\
BLiMP & Zero-shot minimal-pair preference & Morphology, syntax, syntax/semantics, negative polarity items, agreement, islands, binding, argument structure & Average accuracy; higher is better \\
BLiMP Supplement & Zero-shot minimal-pair preference & Additional targeted phenomena including hypernymy, question-answer congruence, subject-auxiliary inversion, and turn taking & Average accuracy; higher is better \\
Entity Tracking & Zero-shot controlled reference task & Whether the model tracks discourse entities through regular, ambiguous-reference, and move-contents conditions across increasing operation counts & Average accuracy; higher is better \\
COMPS & Zero-shot compositional generalization & Compositional interpretation, including nonce/wug-style conditions and distractor variants & Average accuracy; higher is better \\
Reading & Psycholinguistic prediction & Whether model surprisal aligns with human reading behavior in eye-tracking and self-paced reading data & Official reading scores; higher is better \\
\bottomrule
\end{tabular}
}
\caption{Evaluation metrics reported in this paper. The official BabyLM strict zero-shot pipeline scores the available benchmark files in the local evaluation bundle. Accuracy-style tasks compare model preferences over controlled alternatives; reading evaluates alignment between model-derived surprisal and human reading measurements.}
\label{tab:evaluation-benchmarks}
\end{table*}
